\newpage
\begingroup\small\setlength{\emergencystretch}{2em}
\section*{Final DeBERTa model: method and measured results}
\textbf{Status.} The pure neural model scored 0.92433 publicly; the fixed neural/NB-style rank blend scored 0.92164. Both improved on the earlier NB/style blend (0.88628), but the locally better rank blend was worse publicly. Private scores remain unknown.

\subsection*{Hypothesis and experimental protocol}
After observing the validation/public gap, we tested whether contextual pretrained representations could complement local n-grams and stylistic statistics. We used Microsoft's general-purpose \href{https://huggingface.co/microsoft/deberta-v3-small}{DeBERTa-v3-small}, following \href{https://arxiv.org/abs/2111.09543}{DeBERTaV3 (He et al.)}. The exact weight revision is recorded in \texttt{neural\_results/pretrained\_source.json}. Pretraining is external general-language training; the only supervised labels in our fine-tuning are the supplied competition labels. All training and prediction run locally, with no external inference API or hidden labels.

Training uses the same 8,000 training IDs and evaluates the same 2,000 validation IDs. We normalize whitespace, truncate to 256 tokens, freeze pretrained embeddings, and fine-tune the encoder and new classification head. Settings: AdamW, learning rate $3\times10^{-5}$, weight decay 0.01, batch 16, accumulation 2, 10\% linear warmup followed by linear decay, gradient clipping at 1, seed 42, bfloat16 on Apple MPS, two epochs. We retain the better validation checkpoint. Validation labels select the epoch but never enter gradient updates. Test labels are unavailable.

\subsection*{Measured local results}
epoch 1: 0.987573; epoch 2: 0.989423. Selected epoch: 2. Validation-run time (including evaluation and checkpointing): 21.2 minutes.


\par\noindent
\begin{tabularx}{\linewidth}{@{}Yr@{}}\toprule
Model & Fixed-validation AUROC\\\midrule
submitted\_robust\_blend & 0.988136 \\
deberta & 0.989423 \\
deberta\_robust\_rank50 & 0.994594 \\
\bottomrule\end{tabularx}\par

We also test one fixed 50:50 mean-percentile-rank blend of the neural scores and the previously submitted NB/style scores; there is no weight sweep. Ranks are computed separately within validation and test predictions without labels. These ensemble outputs are ranking scores, not calibrated probabilities. Checkpoint and candidate selection reuse the holdout; all comparisons are exploratory.

Selected-checkpoint training AUROC is 0.995930; validation AUROC is 0.989423. Normalizing Unicode, quotation marks and whitespace in both classes gives 0.989299; also removing punctuation gives 0.980220. These are synthetic stress tests, not measured hidden-test conditions.

\subsection*{Limitations and submission procedure}
About 50.5\% of validation texts exceed 256 tokens, so truncation discards potentially useful evidence. A 384-token limit would truncate 8.4\%, but the later combined robustness experiment changes several factors and does not isolate context length. Neural softmax outputs are uncalibrated. Epoch selection and repeated reuse of validation may inflate apparent gains. The observed public scores remain below validation scores, showing that the distribution gap persists.

The final neural candidate is retrained from the pinned pretrained weights on all 10,000 supplied labeled rows using the chosen epoch count and training settings. Its file is \path{artifacts/submission_deberta.csv}; the fixed ensemble is \path{artifacts/submission_deberta_robust_rank50.csv}. Prediction hashes and input hashes are recorded separately, so earlier leaderboard scores cannot transfer to new predictions. See \texttt{NEURAL\_EXPERIMENT.md} for reproduction and limitations. Full-data refit and test inference took 30.3 minutes. The later robustness continuation scored 0.91085 and was rejected; the original pure model remains the final selection.
\endgroup
\newpage
\begingroup\small\setlength{\emergencystretch}{2em}
\section*{Neural candidate: subgroup and error diagnostics}
The same training-derived word-count boundaries (175 and 240 words) define short, medium and long validation groups. AUROC evaluates ranking; the neural classifier's default logit-zero rule corresponds to an uncalibrated softmax value of 0.5.
\par\noindent
\begin{tabularx}{\linewidth}{@{}Yrrr@{}}\toprule
Length group & Count & Neural AUROC & Submitted blend\\\midrule
Long & 646 & 0.991526 & 0.991192 \\
Medium & 643 & 0.987055 & 0.987089 \\
Short & 711 & 0.990983 & 0.986215 \\
\bottomrule\end{tabularx}\par
At the default threshold, the neural model produces 204 false positives and 18 false negatives (88.9\% accuracy). Its high AUROC does not imply a good operating threshold or calibrated probabilities. No validation-optimized threshold is presented as an unbiased test result.

\subsection*{Five false positives and five false negatives}
These are the most confident errors of each type at the default threshold. They are selected examples, not a representative prevalence estimate. Passage inspection supports the following hypotheses; it does not establish causal feature attribution or that supplied labels are incorrect.
\par\noindent
\begin{tabularx}{\linewidth}{@{}rrrY@{}}\toprule
ID & Type & Softmax AI & Possible reason from inspected passage\\\midrule
3326 & FP & 0.999685 & Fluent, conventional review structure resembles generated exposition; surface polish is an ambiguous authorship cue.\\
8238 & FP & 0.999608 & Regular academic framing and repeated technical categories may resemble generated abstracts. The first 256 tokens omit later detail.\\
1761 & FP & 0.999542 & Formulaic research-summary language resembles AI academic prose despite concrete method details.\\
3364 & FP & 0.999512 & Structured instructional prose and repeated terms may be mistaken for generated explanations.\\
1585 & FP & 0.999412 & Predictable abstract structure may dominate the decision; input truncation can hide later quantitative details.\\
2209 & FN & 0.008062 & Recipe conventions are shared by both classes; ingredient-heavy opening text supplies weak authorship evidence.\\
7018 & FN & 0.010987 & Technical surface detail resembles a human abstract; authorship classification does not verify the scientific claims.\\
7837 & FN & 0.013223 & Colloquial imperfections and first-person review style resemble human user comments.\\
424 & FN & 0.016915 & Recipe vocabulary and informal wording resemble human recipes, while truncation removes much of the later sequence.\\
7995 & FN & 0.018833 & Reference-style detail and irregular presentation resemble human source text; the model does not verify biographical accuracy.\\
\bottomrule\end{tabularx}\par
Human research abstracts and structured explanations account for several confident false positives. Recipes, informal reviews and detailed biographies still produce false negatives. These genres occur in both classes. A next controlled experiment could compare 256 versus 384 tokens while retaining the same training/validation IDs; this would test whether later evidence helps rather than assuming truncation caused these errors. More representative held-out domains and probability calibration deserve separate evaluation.
\endgroup
